\documentclass[10pt,a4paper]{article}
\usepackage{amsmath,amssymb,amsthm,graphicx}
\usepackage{bm,enumitem,textcomp,algorithm,xspace,mathtools,subfig,thmtools,thm-restate}
\usepackage[noend]{algpseudocode}
\usepackage[classfont=bold]{complexity}
\usepackage[svgnames]{xcolor}
\usepackage{hyperref}
\usepackage[capitalise,noabbrev,nameinlink]{cleveref}
\graphicspath{{items/}}
\theoremstyle{plain}
\newtheorem{theorem}{Theorem}[section]
\newtheorem{lemma}[theorem]{Lemma}
\newtheorem{corollary}[theorem]{Corollary}
\newtheorem{proposition}[theorem]{Proposition}
\newtheorem{claim}[theorem]{Claim}

\theoremstyle{remark}
\newtheorem{remark}[theorem]{Remark}
\newtheorem{example}[theorem]{Example}

\theoremstyle{definition}
\newtheorem{definition}[theorem]{Definition}
\newcommand{\pa}[1]{}
\newcommand{\kr}[1]{}
%%%
%%% Generic math macros
%%%

% equation with def above
\newcommand\defeq{\stackrel{\mathclap{\normalfont\mbox{\scriptsize{def}}}}{=}}

\DeclareMathOperator*{\supp}{supp}
\DeclareMathOperator*{\conn}{conn}
\DeclareMathOperator{\Char}{char}
\newcommand{\eps}{\varepsilon}

\newcommand{\eqcomma}{\enspace ,}
\newcommand{\eqperiod}{\enspace .}
\newcommand{\disjointunion}{\overset{.}{\cup}}
\newcommand{\expec}[2][]{\mathbb{E}_{#1}[#2]}
\newcommand{\prob}[2][]{\Pr_{#1}[#2]}
\newcommand{\Prob}[2][]{\Pr_{#1} \bigl[ #2 \bigr]}
\providecommand{\abs}[1]{\lvert#1\rvert}
\providecommand{\Abs}[1]{\bigl\lvert#1\bigr\rvert}
\providecommand{\ABS}[1]{\left\lvert#1\right\rvert}
\newcommand{\set}[1]{\{#1\}}
\newcommand{\setdescr}[2]{\{\,#1 \mid #2\,\}}
\newcommand{\floor}[1]{\lfloor #1 \rfloor}
\newcommand{\ceil}[1]{\lceil #1 \rceil}
\newcommand\restrict[2]{{% we make the whole thing an ordinary symbol
  \left.\kern-\nulldelimiterspace % automatically resize the bar with \right
  #1 % the function
  \vphantom{\big|} % pretend it's a little taller at normal size
  \right|_{#2} % this is the delimiter
  }}
% indicator
\usepackage{bbm}
\newcommand{\ind}[1]{\mathbbm{1}{\set{#1}}}

% cal of everything
\newcommand{\calA}{\mathcal{A}}
\newcommand{\calB}{\mathcal{B}}
\newcommand{\calC}{\mathcal{C}}
\newcommand{\calD}{\mathcal{D}}
\newcommand{\calE}{\mathcal{E}}
\newcommand{\calF}{\mathcal{F}}
\newcommand{\calG}{\mathcal{G}}
\newcommand{\calH}{\mathcal{H}}
\newcommand{\calI}{\mathcal{I}}
\newcommand{\calJ}{\mathcal{J}}
\newcommand{\calK}{\mathcal{K}}
\newcommand{\calL}{\mathcal{L}}
\newcommand{\calM}{\mathcal{M}}
\newcommand{\calN}{\mathcal{N}}
\newcommand{\calO}{\mathcal{O}}
\newcommand{\calP}{\mathcal{P}}
\newcommand{\calQ}{\mathcal{Q}}
\newcommand{\calR}{\mathcal{R}}
\newcommand{\calS}{\mathcal{S}}
\newcommand{\calT}{\mathcal{T}}
\newcommand{\calU}{\mathcal{U}}
\newcommand{\calV}{\mathcal{V}}
\newcommand{\calW}{\mathcal{W}}
\newcommand{\calX}{\mathcal{X}}
\newcommand{\calY}{\mathcal{Y}}
\newcommand{\calZ}{\mathcal{Z}}


% R, N, Z, F etc
\renewcommand{\C}{\mathbb{C}}
\renewcommand{\E}{\mathbb{E}}
\newcommand{\PE}{\mathbb{\widetilde{E}}}
\renewcommand{\R}{\mathbb{R}}
\newcommand{\N}{\mathbb{N}}
\newcommand{\Z}{\mathbb{Z}}
\newcommand{\F}{\mathbb{F}}

\newcommand{\trace}{\mathsf{tr}}
\newcommand{\Var}{\mathsf{Var}}


\DeclareMathOperator{\Deg}{\sf{Deg}}
\DeclareMathOperator{\Size}{\sf{Size}}
\newcommand{\RefuteDeg}[2]{\Deg(#2 \vdash_{\text{\tiny #1}} \, \perp)}
\newcommand{\RefuteSize}[2]{\Size(#2 \vdash_{\text{\tiny #1}} \, \perp)}

% pseudo-distribution
\newcommand{\pedef}{\tilde{\mathbb{E}}}
\newcommand{\pe}[2][]{\pedef_{#1}[#2]}

% monotone anything
\newcommand{\mon}{{\mathsf{mon}}}

% circuit size
\DeclareMathOperator{\csize}{\calC}
\DeclareMathOperator{\cmonsize}{\csize_{\mathsf{mon}}}

% Misc
\DeclareMathOperator{\hw}{\texttt{hw}}

% out and in for the g fns
\newcommand{\gout}{g^\mathrm{out}}
\newcommand{\ginone}{g^\mathrm{in_1}}
\newcommand{\gintwo}{g^\mathrm{in_2}}

% Notation related to the expander graphs
\DeclareMathOperator{\idx}{\texttt{ind}}
\DeclareMathOperator{\neigh}{\textbf{Neigh}}
  

% subformula
\newcommand{\subformula}{$ is a variable substitution of $}

% the selectors
\DeclareMathOperator{\sel}{\mathsf{Sel}}

% the circuit formula
\newcommand{\vars}[1]{\mathsf{#1}}
\DeclareMathOperator{\circuit}{Circuit}
\DeclareMathOperator{\outwire}{\vars{Out}}
\DeclareMathOperator{\inwire}{\vars{In}}
\DeclareMathOperator{\isneg}{\vars{IsNeg}}
\DeclareMathOperator{\isand}{\vars{IsAnd}}
\DeclareMathOperator{\isor}{\vars{IsOr}}
\DeclareMathOperator{\isfromconst}{\vars{IsFromConst}}
\DeclareMathOperator{\isfromvar}{\vars{IsFromInput}}
\DeclareMathOperator{\isfromgate}{\vars{IsFromGate}}
\DeclareMathOperator{\constval}{\vars{ConstantValue}}
\DeclareMathOperator{\isvar}{\vars{IsInput}}
\DeclareMathOperator{\isgate}{\vars{IsGate}}

% axioms for cref
\crefname{axiom}{Axiom}{Axioms}
\Crefname{axiom}{Axiom}{Axioms}

\crefname{axioms}{Axioms}{Axioms}
\Crefname{axioms}{Axioms}{Axioms}
% old notation -- might need if we create an appendix:
% cnf <-> poly encoding
\DeclareMathOperator{\isvarless}{\vars{IsInput^{\le}}}
\DeclareMathOperator{\isgateless}{\vars{IsGate^{\le}}}
\title{Linear-degree Sum-of-Squares refutation upper bound for false circuit-size claims}
\author{Per Austrin, Kilian Risse}
\date{doi:10.4230/LIPIcs.CCC.2023.31; CC BY4.0}
\begin{document}
\maketitle
\setcounter{section}{1}
\noindent\textit{Editorial note.} Original authors: Per Austrin and Kilian Risse. arXiv:2311.12994v1 (2023-11-21), https://arxiv.org/abs/2311.12994v1. Creative Commons Attribution4.0, https://creativecommons.org/licenses/by/4.0/.\par\medskip
\noindent\textit{Editorial note.} Selected complete native author main theorem/proof/context, original circuit diagrams and bibliography in standalone layout. Input files explicitly flattened into source spans, unused frontmatter/font packages omitted. No mathematical changes or new proofs. Original CC BY4.0 license retained.\par\medskip
\noindent\textit{Allowed standard prerequisites (editorial).}\begin{itemize}
\item Sum-of-Squares derivations, Boolean-cube polynomial identities and circuit polynomial encoding. Full author preliminaries and all Circuit axioms retained.
\end{itemize}\medskip
\section{Preliminaries}
\label{sec:preliminaries}

All logarithms are in base $2$.  For integers $n \ge 1$ we write
$[n] = \set{1, 2, \ldots, n}$ and for a set $U$ we denote the power
set of $U$ by $2^{U}$. Further, for a set $V \subseteq U$ we let
$\overline V$ be the complement of $V$ with respect to $U$, that is,
$\overline V = U \setminus V$. We write
$\binom{[n]}{\ell} \subseteq \set{0,1}^n$ for the set of binary
strings with Hamming weight $\ell$. For a string
$\alpha \in \set{0,1}^n$ we let
$|\alpha| = \sum_{i \in [n]} \alpha_i$.

We sometimes want to supress dependencies on constants and write
$f(n, \eps) \in O_\eps\big(g(n, \eps)\big)$, respectively
$f(n, \eps) \in \Omega_\eps\big(g(n, \eps)\big)$, to mean that there
exists a function $c(\eps) > 0$ such that there is an $n_0$ and for all
$n \ge n_0$ it holds that $f(n, \eps) \le c(\eps) \cdot g(n, \eps)$,
respectively $f(n, \eps) \ge c(\eps)\cdot g(n, \eps)$.

\begin{definition}\label{def:explicit}
  A sequence of bipartite graphs $\set{G_n=(U_n, V_n, E_n)}_{n\in \N}$
  with $\deg(u) =d$ for all $u \in U_n$ is \emph{explicit} if there is
  an algorithm that given $(n, u, j)$, where $n \in \N, u \in U_n$ and
  $j\in [d]$, computes the $j$th neighbor of vertex $u$ in the graph
  $G_n$ in time $\poly(\log n + \log |U| + \log d)$.
\end{definition}

From now on it is understood that whenever we talk about an explicit
graph we actually mean to say that there is a sequence of explicit
graphs with above properties.

For a graph $G = (V, E)$ and $W \subseteq V$ we denote by
\[
  N(W) = \setdescr{u \in V \setminus W}{(w,u) \in E \text{ for some } w
    \in W}
\]
the set of neighbors of $W$.

\begin{definition}
  A bipartite graph $G = (U, V, E)$ is an $(r, d, c)$-expander if every
  vertex $u \in U$ has degree $\deg(u) = d$ and every set
  $W \subseteq U$ of size $|W| \le r$ satisfies $|N(W)| \ge c \cdot |W|$.
\end{definition}

A key ingredient in our proofs is the following result on the
existence of strong explicit expanders.

\begin{theorem}[\cite{GUV09Unbalanced}]\label{thm:expander}
  For all constants $\gamma > 0$,
  every $M \in \N$, $r \le M$,
  and $\eps > 0$,
  there is an
  $N \le d^2 \cdot r^{1 + \gamma}$ and an
  explicit
  $(r, d, (1- \eps)d)$-expander
  $G=(U, V, E)$,
  with $|U| = M$,
  $|V| = N$, and
  $d = O\big( ((\log M)(\log r) / \eps )^{1+1/\gamma}\big)$.
\end{theorem}

\kr{Igor pointed this out:
  \url{https://eccc.weizmann.ac.il/report/2018/159/}. Quickly checking
  the parameters I believe it does not give anything better. Need to
  check more carefully}

For our purposes it is more relevant to compute the neighbor relation
$\neigh(u,v)$ indicating whether $(u, v) \in E$ rather than the
neighbor function as in \cref{def:explicit}, but this is an immediate
consequence of being able to compute the neighbor function.

\begin{claim}
  \label{claim:explicit neigh}
  If $G = (U, V, E)$ is explicit then the neighbor relation
  $\neigh: U \times V \rightarrow \{0,1\}$ is computable by a circuit
  of size
  $d \cdot \big( \poly(\log n + \log |U| + \log d) + 2 \log |V| + 1
  \big)$.
\end{claim}

% \pa{Added this fact, check that I was not too optimistic in claimed
  % size bound. Maybe add proof here.} \kr{checked.}

% Recall that this graph is explicit and hence there are
% circuits to compute the neighbor function
% $\Gamma: U \times [k] \rightarrow Z$.

% Each circuit $C_1, \ldots, C_z$ consists of $k$ circuits $\calD =
% \set{D_1, \ldots, D_k}$ that compute $\Gamma$ with the second input
% fixed to their corresponding index. On top of these circuits from
% $\calD$ we add a small circuit  $E_i$, which computes the following.
% Each circuit $D \in \calD$ outputs
% some index $j \in [z]$. The circuit $E_i$ outputs $1$ if and only if
% one of the $D \in \calD$ outputs the index $i$. The circuit $E_i$ is
% fairly small: for each circuit $D$ we require $2 \log z$ gates to
% check whether $D$ output $i$, and an additional $k$ gates to \emph{or}
% these outputs together. So in total each $E_i$ consists of
% $k (2\log z + 1)$ many gates.
% As the graph $G$ is explicit, we know that there is a circuit of size
% $\poly(n + \log k)$ computing $\Gamma$ and hence the size of all
% circuits in $\calD$ is bounded by $c_\alpha \cdot k \cdot n^c$, for
% some constant $c$, and $c_\alpha$ depending on $\alpha$. Hence we
% conclude that
% \begin{align}\label{eq:bound-a}
%  a = z \cdot \big(k (2\log z + 1) + c_\alpha \cdot k \cdot n^c\big) \eqperiod
% \end{align}


A slice function is a Boolean function $f$ such that there is a
$\ell \in [n]$ with $f(\alpha) = 0$ whenever $|\alpha| < \ell$, and
$f(\alpha) = 1$ whenever $|\alpha| > \ell$.  Note that all slice
functions are monotone.

The circuit complexity $\csize(f)$ of a Boolean function $f$ is the
size of the smallest circuit over the basis $\vee, \wedge$, and
$\lnot$ (with fan-in $2$). Similarly the monotone circuit
complexity $\cmonsize(f)$ of a monotone Boolean function $f$
is the size of the smallest circuit over the basis $\vee$, and
$\wedge$. We have the following useful inequality between these
measures.

\begin{lemma}[\cite{B82}]\label{lem:slice}
  If $g$ is any slice function on $n$ bits, then
  $\cmonsize(g) \le 2 \csize(g) + O(n^2 \log n)$.
\end{lemma}

Finally we also rely on the following simple claim.

\begin{claim}
  \label{claim:0 mod bool}
  Let $p: \R^n \rightarrow \R$ be a degree $d$ polynomial such that
  $p(x) = 0$ for all $x \in \{0,1\}^n$.  Then $p$ can be written as
  \[
    p(x) = \sum_{i \in [n]} q_i(x) \cdot (x_i^2-x_i)
  \]
  where each term in the sum has degree at most $d$.
\end{claim}

\begin{proof}[Proof sketch]
  We take the polynomial $p$ and multilinearize it, using the
  appropriate polynomial $x_i^2 - x_i$. Eventually we are left with a
  sum of polynomials of the form $q_i(x)\cdot(x_i^2 - x_i)$ and a
  multilinear polynomial $\tilde p(x)$ which is $0$ on all Boolean
  inputs.  As multilinear polynomials are a basis for Boolean
  functions this implies that $\tilde p(x)$ is equal to the $0$
  polynomial and hence the claim follows.
\end{proof}

\subsection{Sum of Squares}
\label{sec:known-SoS-results}

Let $\calP = \set{p_1 = 0, \ldots, p_m = 0}$ be a system of polynomial
equations over the set of variables
$X = \set{x_1, \ldots, x_n, \bar{x}_1, \ldots, \bar{x}_n}$.  Each
$p_i$ is called an \emph{axiom}, and throughout the paper we always
assume that $\calP$ includes all axioms $x_i^2 - x_i$ and
$\bar{x}_i^2 - \bar{x}_i$, ensuring that the variables are Boolean, as
well as the axioms $1 - x_i - \bar{x}_i$, making sure that the ``bar''
variables are in fact the negation of the ``non-bar'' variables.

%

\begin{definition}[Sum-of-Squares]
  \emph{Sum-of-Squares (SoS)} is a static, semi-algebraic proof
  system. An SoS proof of $f \ge 0$ from $\calP$ is a sequence of
  polynomials $\pi = (t_1, \ldots, t_m; s_1, \ldots, s_a)$ such that
  \begin{align*}
    \sum_{i \in [m]} t_i p_i +
    \sum_{i \in [a]} s^2_i = f \eqperiod %\label{eq:sos-def}
  \end{align*}
  The \emph{degree} of a proof $\pi$ is
  $\Deg(\pi) = \max \set{ \max_{i \in [m]} \deg(t_i) + \deg(p_i),
    \max_{i \in [a]} 2 \deg(s_i) }$, an \emph{SoS refutation of
    $\calP$} is an SoS proof of $-1 \ge 0$ from $\calP$, and the
  \emph{SoS degree to refute $\calP$} is the minimum degree of any SoS
  refutation of $\calP$: if we let $\pi$ range over all SoS
  refutations of $\calP$, we can write
  $\RefuteDeg{SoS}{\calP} = \min_\pi \Deg(\pi)$.  The \emph{size of an
    SoS refutation $\pi$}, $\Size(\pi)$, is the sum of the number of
  monomials in each polynomial in $\pi$ and the \emph{size of refuting
    $\calP$} is the minimum size over all refutations
  $\RefuteSize{SoS}{\calP} = \min_{\pi} \Size(\pi)$.
\end{definition}


Let us recall some well-known results about SoS. Given a bipartite
graph $G = (U, V, E)$, and $b \in \set{0,1}^{|U|}$ we denote by
$\Phi(G,b)$ the following XOR-CSP instance defined over $G$: for each
$v \in V$ there is a Boolean variable $x_v$, and for every vertex
$u \in U$ there is a constraint $\oplus_{v \in N(u)} x_v = b_u$.
We encode this in the obvious way as a system of polynomial equations: 
\begin{align*}
  \big\{
  \prod_{v \in N(u)} (1 - 2 \cdot x_v) = 1 - 2 \cdot b_u
  \mid u \in U
  \big\} \eqcomma
\end{align*}
along with the Boolean axioms and the negation axioms for the $x$
variables. The first theorem we need to recall is the classic lower
bounds for XOR-CSPs by Grigoriev.

\begin{theorem}[\cite{gri01xor}]\label{thm:sos-xor}
  For $n \in \N$, all $k = k(n)$ and $r = r(n)$ the following holds.
  Let $G = (U, V, E)$ be an $(r, k, 2)$-expander with $|V| = n$. Then
  for every $b \in \set{0,1}^{|U|}$ SoS requires degree $\Omega(r)$ to
  refute the claim that there is a satisfying assignment to
  $\Phi(G, b)$.
\end{theorem}

We also need to recall the size-degree tradeoff by Atserias and
Hakoniemi.

\begin{theorem}[\cite{AH18SosTradeoff}]\label{thm:tradeoff}
  Let $\calP$ be a system of polynomial equations over $n$ Boolean
  variables and degree at most $k$. If $d$ is the minimum degree
  SoS requires to refute $\calP$, then the minimum size of an SoS
  refutation of $\calP$ is at least $\exp(\Omega((d-k)^2/n))$.
\end{theorem}


\subsection{Restrictions}
\label{sec:restrictions}

Let $\calP = \set{p_1 = 0, \ldots, p_m = 0}$ be a system of polynomial
equations over the set of Boolean variables
$X = \set{x_1, \ldots, x_n, \bar{x}_1, \ldots, \bar{x}_n}$.
%
For a map
$\rho: \set{x_1, \ldots, x_n} \to \set{0,1, x_1, \ldots, x_n,
  \bar{x}_1, \ldots, \bar{x}_n}$ denote by $\restrict{\calP}{\rho}$
the system of polynomial equations $\calP$ restricted by $\rho$, i.e.,
\begin{align*}
\restrict{\calP}{\rho} = \set{ &p_1(\rho(x_1), \ldots, \rho(x_n)) = 0,\\
  &p_2(\rho(x_1), \ldots, \rho(x_n)) = 0,\\ &\vdots\\ &p_m(\rho(x_1),
  \ldots, \rho(x_n)) = 0 }\eqcomma
\end{align*}
where it is understood that $\rho(\bar x_i) = \overline{\rho(x_i)}$,
with the convention $\bar{\bar{x}}_i = x_i$, $\bar 0 = 1$ and vice
versa. Throughout the paper all our restrictions set the bar variables
to the negation of the non-bar variables. As such it makes sense to
treat the pair of variables $(x_i, \bar x_i)$ as one variable and we
say that $\calP$ has $n$ \emph{unset} variables.

\begin{definition}[Variable Substitution]
  We say that a system of polynomial equations $\calP'$ is a
  \emph{variable substitution of $\calP$} if there is a map
  $\rho: \set{x_1, \ldots, x_n} \to \set{0,1, x_1, \ldots, x_n,
    \bar{x}_1, \ldots, \bar{x}_n}$ such that
  $\calP' = \restrict{\calP}{\rho}$, where we ignore polynomial
  equations of the form $0=0$.
\end{definition}

The following well-known lemma states that a system of polynomial
equations $\calP$ is at least as hard as any of its variable substitutions.
\begin{lemma}\label{lem:subformula}
  Let $\calP, \calP'$ be systems of polynomial equations such that
  $\calP' \subformula \calP$. Then,
  \begin{enumerate}[label=$(\roman*)$]
  \item $\RefuteDeg{SoS}{\calP} \ge \RefuteDeg{SoS}{\calP'}$, and
  \item $\RefuteSize{SoS}{\calP} \ge \RefuteSize{SoS}{\calP'}$.
  \end{enumerate}
\end{lemma}
The lemma is easy to verify by considering an SoS refutation of
$\calP$ and hitting it with the appropriate variable substitution. The
restricted proof is now a refutation of $\calP'$ and it can be seen
that the degree/size of the restricted refutation is at most the
degree/size of the original refutation.

We also consider more general substitutions.

\begin{definition}[Polynomial Substitution]
  Functions
  $\rho: \set{x_1, \ldots, x_n} \rightarrow \R[x]_{\le k}$ that map
  variables to polynomials of degree at most $k$ are called
  \emph{polynomial substitutions}.
\end{definition}

For polynomial substitutions we have the following well-known lemma.

\begin{lemma}
  \label{lem:sos degree polynomial substitution}
  Let $\calP$ be a system of polynomial equations and let $\rho$ be a
  polynomial substitution mapping variables to polynomials of degree
  at most $k$. Then,
  $\RefuteDeg{SoS}{\calP} \ge
  \RefuteDeg{SoS}{\restrict{\calP}{\rho}}/k$.
\end{lemma}

This lemma can again be verified by considering a refutation of
$\calP$. Substitute each variable $x_i$ in the proof by
$\rho(x_i)$. This results in a refutation of $\restrict{\calP}{\rho}$,
whose degree is at most a factor $k$ larger than the degree of the
refutation of $\calP$.

\subsection{The Circuit Size Formula}
\label{sec:circuit-encoding}
The formula $\circuit_s(f)$ encodes the claim that the function $f$,
given as a truthtable $f \in \set{0, 1}^{2^n}$, can be computed by a
circuit of size $s$ over $n$ Boolean inputs $x_1, \ldots, x_n$. The
encoding is not essential but for concreteness let us fix one encoding
of this claim. We deviate from the encoding used by Razborov
\cite{Razborov98, Razborov04ResolutionLowerBoundsPM} and do not
present the formula as a propositional formula but rather as a system
of polynomial equations. In order to encode below constraints as a
constant width CNF formula, as done by Razborov, one needs to
introduce extension variables. Despite this difference it is not
difficult to see that our lower bound also works against the CNF
encoding. In \cref{sec:encoding} we directly show that a low degree
SoS refutation of the CNF encoding gives rise to a low degree SoS
refutation of the encoding used in this paper (see
\cref{prop:cnf-poly}). Thus a lower bound for our encoding implies a
lower bound for the CNF encoding. As the presentation is simpler in
the polynomial encoding, we present it as follows.

We also need to define the monotone version of $\circuit_s(f)$ denoted
by $\circuit^{\mathsf{mon}}_s(f)$. The later is a restriction of the
former with the $\isneg(v)$ (see below) variable, for all $v \in [s]$,
set to $0$. This forces the circuit to only contain $\land$ and $\lor$
gates, i.e., the circuit is monotone.

All variables introduced in the following are Boolean variables and we
implicitly add the Boolean axiom $y(1 - y) = 0$ for each variable $y$
and further implicitly introduce the ``bar variable'' $\bar{y}$ along
with the negation axiom $y = 1 - \bar{y}$ (and the corresponding
Boolean axiom) ensuring that $\bar{y}$ is always the negation of $y$.

Let us first describe the \emph{structure variables} which are used to
describe the circuit that supposedly computes the function $f$.

We view the $s$ gates as being indexed from $1$ to $s$ in topological
order with gate $s$ being the output.
For each gate $v \in [s]$ there are three variables $\isneg(v),
\isor(v), \isand(v)$ indicating the operation computed at
$v$. Similarly for a gate $v \in [s]$ and a wire $a \in \set{1,2}$ we
have variables $\isfromconst(v, a), \isfromvar(v, a), \isfromgate(v, a)$
indicating whether the input wire $a$ of $v$ is connected to a
constant, a variable or a gate.

Further, we have the variables $\constval(v, a)$, $\isvar(v, a, i)$
and $\isgate(v, a, u)$, for $a \in \set{1,2}$, $i \in [n]$ and
$u < v$, specifying the constant value, input $x_i$ or gate $u$, the
input wire $a$ of $v$ is connected to (assuming $a$ is connected to
the corresponding kind).

The second set of variables are the \emph{evaluation variables}, which
describe what value is computed at each $v$ on input
$\alpha = \alpha_1, \ldots, \alpha_n$ (i.e., we have $x_i = \alpha_i$).

For each gate $v \in [s]$ and assignment $\alpha \in \set{0,1}^n$ we
have a Boolean variable $\outwire_\alpha(v)$ indicating the value
computed at gate $v$ on input $\alpha$. The Boolean variable
$\inwire_{\alpha}(v, a)$ indicates the value brought to the vertex
$v \in [s]$ on wire $a \in \set{1,2}$ on input $\alpha$.

Note that there is a total of
$3s + 6s + 2s + 2sn + 2{s \choose 2} = \Theta(s^2+sn)$ structure
variables, and a total of $3 s 2^n$ evaluation variables, for a total
of $\Theta(s^2 + s 2^n)$ variables in $\circuit_s(f)$.

The formula consists of the following axioms. For the sake of
readability we omit some universal quantifiers: the variable
$a \in \set{1,2}$ in
\cref{eq:ax-from,eq:ax-in-gt,eq:ax-in-gt2,eq:prod-zero-in,eq:prod-zero-gate,eq:inwire-const,eq:inwire-var,eq:inwire-gate}
as well as the variable $\alpha \in \set{0,1}^n$ in
\cref{eq:inwire-const,eq:inwire-var,eq:inwire-gate,eq:ax-neg,eq:ax-or,eq:ax-and,eq:ax-correct}
are implicitly universally quantified.

Let us first describe
the axioms on the structure of the circuit. In the following
section we refer to this set of axioms as the \emph{structure
  axioms}. The first axioms ensure that every wire is connected to a
single kind
\begin{align}
  \label{eq:ax-from}
  \isfromconst(v, a) +
  \isfromvar(v, a) +
  \isfromgate(v, a) = 1 \quad
  \forall \, v \in [s] \eqcomma
\end{align}
and similarly the next axioms make sure that each gate is of
precisely one kind
\begin{align}
  \label{eq:ax-fn}
  \isneg(v) + \isor(v) + \isand(v) = 1
  \quad \forall \, v \in [s] \eqperiod
\end{align}
The final structure axioms ensure that the variables, which indicate
to what input or gate a fixed wire is connected to, always sum to one
(except for gate $1$ which cannot have any inputs from other gates)
\begin{align}
  \label{eq:ax-in-gt}
  \sum_{i = 1}^{n}\isvar(v,a,i) &= 1
  \quad \forall v \in [s], \text{ and}\\
  \label{eq:ax-in-gt2}
  \sum_{u = 1}^{v-1}\isgate(v,a,u) &= 1
  \quad \forall v \in [s] \setminus\{1\} \eqperiod
\end{align}
We further strengthen our encoding by adding the axioms
\begin{align}
  \label{eq:prod-zero-in}
  \isvar(v,a,i)\isvar(v,a,j) &= 0
  \quad\forall v \in [s],\, i < j \in [n], \text{ and}\\
  \label{eq:prod-zero-gate}
  \isgate(v,a,u)\isgate(v,a,u') &= 0
  \quad\forall u < u' < v \in [s] \eqperiod
\end{align}
Note that \cref{eq:prod-zero-in,eq:prod-zero-gate} are implied by
\cref{eq:ax-in-gt,eq:ax-in-gt2}. We add these axioms in order to argue
that a short refutation of the CNF encoding of this principle leads to
a short refutation of the present encoding.

The second group of axioms are the \emph{evaluation axioms} and they
ensure that the evaluation variables indeed compute the intended
values. We start by making sure that the wires carry the value
intended by the structure axioms. If a wire is connected to a
constant, then the evaluation variable associated with that wire
should always be equal to the constant
\begin{align}
  \label{eq:inwire-const}
  \isfromconst(v,a) \cdot
  \big( \inwire_{\alpha}(v,a) - \constval(v, a) \big) = 0 \eqcomma
\end{align}
and similarly in case if a wire is connected to an input or a gate
\begin{align}
  \label{eq:inwire-var}
  \isfromvar(v, a) \cdot \isvar(v, a, i) \cdot
  \big( \inwire_\alpha(v, a) - \alpha_i\big) &= 0
  \eqcomma\\
  \label{eq:inwire-gate}
  \isfromgate(v, a) \cdot \isgate(v, a, u) \cdot
  \big(\inwire_\alpha(v, a) - \outwire_\alpha(u)\big) &= 0 \eqperiod
\end{align}
The final set of evaluation axioms makes sure that the output
evaluation variable of a gate is correctly related to the input
evaluation variables:
\begin{align}
  \label{eq:ax-neg}
  \isneg(v) \cdot
  \outwire_{\alpha}(v)
  &=
    \isneg(v) \cdot
    \overline{\inwire_{\alpha}(v, 1)}\eqcomma\\
  \label{eq:ax-or}
  \isor(v) \cdot
  \outwire_{\alpha}(v)
  &=
    \isor(v)
    \cdot \big(
    1 -
    \overline{\inwire_{\alpha}(v, 1)} \cdot
    \overline{\inwire_{\alpha}(v, 2)}
    \big) \eqcomma\\
  \label{eq:ax-and}
  \isand(v) \cdot
  \outwire_{\alpha}(v)
  &=
    \isand(v) \cdot
    \inwire_{\alpha}(v, 1) \cdot
    \inwire_{\alpha}(v, 2) \eqperiod
\end{align}

Last but not least we have the axioms that ensure that the circuit
outputs the function specified by the truthtable
\begin{align}
  \label{eq:ax-correct}
  \outwire_{\alpha}(s) &= f(\alpha) \eqperiod
\end{align}

\begin{restatable}{theorem}{MainDegree}
  \label{thm:main-degree}
  For all $\eps > 0$ there is a $d = d(\eps)$ such that the following
  holds.  For $n \in \N$, all $s \ge n^{d}$ and any Boolean function
  $f:\set{0,1}^{n} \rightarrow \set{0,1}$ on $n$ bits, SoS requires
  degree $\Omega_\eps(s^{1-\eps})$ to refute $\circuit_s(f)$.
\end{restatable}

It is worthwhile to point out that the proof of \cref{thm:main-degree}
is not specific to the SoS proof system. In fact we outline a general
reduction that shows that if one has a CSP lower bound of the form of
\cref{thm:sos-xor} that only requires good expansion of the underlying
constraint-variable incidence graph and the proof system is expressive
enough so that one can replace variables by local Boolean functions,
then one obtains strong lower bounds for the $\circuit_s(f)$ formula.

The lower bound of $\Omega_\eps(s^{1-\eps})$ on the degree is
essentially tight: if $f$ does not have a circuit of size $s$ then
there exists an SoS refutation of this statement in degree $O(s)$.

\section[On Encodings of the $\circuit_s(f)$ Tautology]{On Encodings
  of the $\bm{\mathrm{Circuit}_s(f)}$ Tautology}
\label{sec:encoding}

Let us introduce a possible constant width CNF encoding of
$\circuit_s(f)$ as proposed by Razborov \cite{Razborov98}.

The formula is defined over the same variables as introduced in
\cref{sec:circuit-encoding}, but in order to keep the fan-in bounded,
we further introduce the extension variables $\isvarless(v, a, i)$ and
$\isgateless(v, a, u)$ that indicate whether the wire $a$ of $v$ is
connected to a variable in $x_1, \ldots, x_i$, a gate $1, \ldots, u$
respectively.

Let us group the axioms in the same manner as we did in
\cref{sec:circuit-encoding}. First we have the structure axioms. The
first axioms encode that each wire is connected to a single kind
\begin{align}
  &\big(\isfromconst(v, a) \vee \isfromvar(v, a) \vee \isfromgate(v,
    a)\big) \wedge \nonumber\\
  &\neg\big(\isfromconst(v, a) \wedge \isfromvar(v, a)\big) \wedge\nonumber\\
  &\neg\big(\isfromvar(v, a) \wedge \isfromgate(v, a)\big) \wedge\nonumber\\
  &\neg\big(\isfromconst(v, a) \wedge \isfromgate(v, a)\big) \eqperiod
\end{align}
The next set of axioms similarly ensures that each gate computes
precisely one function
\begin{align}
  &\big(\isneg(v) \vee \isor(v) \vee \isand(v)\big) \land\nonumber\\
  &\neg\big(\isneg(v) \wedge \isor(v)\big) \wedge \neg\big(\isor(v) \wedge
    \isand(v)\big) \wedge \neg\big(\isneg(v) \wedge \isand(v)\big) \eqperiod
\end{align}
Last, we need to make sure that each wire is connected to a single
input or a gate.
\begin{align}
  \label{ax:input-cnf}
  &\isvarless(v, a, n) \wedge \bigwedge_{i \neq j}\neg \big(\isvar(v,a,i) \wedge
    \isvar(v,a,j)\big) \wedge\nonumber\\
    &\bigwedge_{i \in [n]}
      \Big(\isvarless(v, a, i) \equiv
    \big(\isvarless(v, a, i-1) \vee \isvar(v, a, i)\big)
    \Big) \eqcomma \nonumber\\
  &\qquad\text{where } \isvarless(v, a, 0) \defeq 0 \eqcomma
\end{align}
and similarly for $v \in [s] \setminus \{1\}$ we have that 
\begin{align}
  \label{ax:gate-cnf}
  &\isgateless(v, a, v-1) \wedge
    \bigwedge_{u < u' < v} \neg\big(\isgate(v,a,u) \wedge \isgate(v,a,u')\big)
    \wedge\nonumber\\
  &\bigwedge_{u \in [v-1]}
    \Big(
    \isgateless(v, a, u) \equiv
    \big(\isgateless(v, a, u-1) \vee \isgate(v, a, u)\big)
    \Big)\eqcomma \nonumber\\
  &\qquad\text{where }\isgateless(v, a, 0) \defeq 0 \eqperiod
\end{align}

Let us take a look at the evaluation axioms. Again, we have axioms
that ensure that the wires carry the values intended by the structure
variables. If a wire is connected to a constant, then the evaluation
variable associated with that wire should be equal to the constant
\begin{align}
  &\isfromconst(v, a)
    \rightarrow
    \big(\inwire_{
    \alpha}(v, a) \equiv \constval(v, a)\big) \eqcomma
\end{align}
and similarly if a wire is connected to an input or a gate
\begin{align}
  &\isfromvar(v, a) \wedge \isvar(v, a, i)
    \rightarrow
    \inwire_{\alpha}(v, a) \equiv \alpha_i \eqcomma\\
  &\isfromgate(v, a) \wedge \isgate(v, a, u)
    \rightarrow
    \inwire_{\alpha}(v, a) \equiv \outwire_{\alpha}(u) \eqperiod
\end{align}
Last we need to make sure that the gates propagate the value they are
supposed to compute.
\begin{align}
  &\isneg(v) \rightarrow
    \big(\outwire_{\alpha}(v) \equiv \neg \inwire_{\alpha}(v, 1)\big)\\
  &\isor(v) \rightarrow
    \big(\outwire_{\alpha}(v) \equiv \inwire_{\alpha}(v, 1) \vee \inwire_{\alpha}(v, 2)\big)\\
  &\isand(v) \rightarrow
    \big(\outwire_{\alpha}(v) \equiv \inwire_{\alpha}(v, 1) \wedge
    \inwire_{\alpha}(v, 2)\big)
    \eqperiod
\end{align}

The final axioms ensure that the correct function is computed
\begin{align}
  \outwire_{\alpha}(s) \equiv f(\alpha) \eqperiod
\end{align}

This formula can be rewritten in the usual manner into a $4$-CNF. Let
us denote this formula by $\circuit^{\mathrm{CNF}}_s(f)$.

\pa{Here the $4$ is because all axioms use at most 4 variables, right?  Can say this.}

\kr{It is a bit annoying: you have to consider it as a CNF and then
  each clause is over at most $4$ vars. The axioms themselves (for
  example \cref{ax:gate-cnf}) may depend on more variables. }

\begin{proposition}
  \label{prop:cnf-poly}
  If there is an SoS refutation of degree $d$ of the CNF formula
  $\circuit^{CNF}_s(f)$, then there is an SoS refutation of degree
  $O(d)$ of the system of polynomials $\circuit_s(f)$ as introduced in
  \cref{sec:circuit-encoding}.
\end{proposition}

The rest of this section is devoted to the proof of
\cref{prop:cnf-poly}.

Observe that for each axiom $p$ from the polynomial encoding
$\circuit_s(f)$, there is a CNF
$F_p \subseteq \circuit_s^{\mathrm{CNF}}(f)$ over the same variables
as $p$ (ignoring the added extension variables $\isvarless(v, a, i)$
and $\isgateless(v, a, u)$) such that $p(\alpha) = 0$ is satisfied by
a Boolean assignment $\alpha$ if and only if $F_p$ is satisfied by
$\alpha$ (where we extend the assignment to the extension variables in
the natural manner).

\pa{``extension variables'' not defined.  When introducing them, say
  that they are extension variables, and remind here what the
  extension variables are.  But I also don't understand what the
  purpose of this paragraph is, why is it needed?}

\kr{added the definition of extension variables and remind the reader
  by adding the names. The purpose of this paragraph is to introduce
  the notation $F_p$.}

Recall that SoS operates on polynomials and we thus need to translate
the CNF into a system of polynomials. We translate a clause
$\vee_{i \in [w]} z_i$ into the polynomial
$\prod_{i \in [w]}(1 - z_i) = 0$.

Observe that almost all axioms $p$ of $\circuit_s(f)$ depend only on a
constant number of variables. From such $p$, using the appropriate
Boolean axioms and negation axioms, we can in constant degree derive
$F_p$.

Let us define a polynomial substitution $\rho$ that gets rid of the
extension variables in the natural manner: the substitution $\rho$
first replace each occurrence of a ``bar'' extension variable
$\overline{\isvarless(v, a, i)}$ or $\overline{\isgateless(v, a, u)}$
by the polynomial $1-\isvarless(v, a, i)$ and the polynomial
$1-\isgateless(v, a, u)$ respectively. Then, $\rho$ replaces each
occurrence of the variable $\isvarless(v, a, i)$ by
$\sum_{j \le i} \isvar(v, a, j)$ and similarly $\isgateless(v, a, u)$
by $\sum_{w \le u} \isgate(v, a, w)$.

Suppose we have a degree $d$ refutation $\pi$ of
$\circuit_s^{\mathrm{CNF}}(f)$. Let us consider $\restrict{\pi}{\rho}$
and $\restrict{\circuit^{\mathrm{CNF}}(f)}{\rho}$. Note that
$\restrict{\pi}{\rho}$ is a degree $d$ SoS refutation of
$\restrict{\circuit^{\mathrm{CNF}}_s(f)}{\rho}$.

We claim that in constant degree the axioms of
$\restrict{\circuit^{\mathrm{CNF}}_s(f)}{\rho}$ can be derived from the
polynomial encoding $\circuit_s(f)$. As previously noted, this holds
for all axioms but the ones that are over a non-constant number of
variables. In other words it just remains to show that we can derive
the substituted \cref{ax:input-cnf,ax:gate-cnf} from
\cref{eq:ax-in-gt,eq:ax-in-gt2,eq:prod-zero-in,eq:prod-zero-gate}.

Let us consider \cref{ax:input-cnf}. With the extension variables
substituted and translated into a system of polynomials the axiom
consists of the following polynomial equations.
\begin{align}
  \label{eq:ax1}
  1 - \sum_{j \in [n]} \isvar(v,a,j) &= 0\\
  \label{eq:ax2}
  \isvar(v,a,i) \cdot \isvar(v,a,j) &= 0,
  \text{ for~}i \neq j\\
  \nonumber
  \left(\sum_{j \le i} \isvar(v,a,j)\right)
  \left(1 - \sum_{j < i} \isvar(v,a,j)\right)\cdot\quad&\\
  \label{eq:ax3}
  \overline{\isvar(v,a,i)} &= 0,
  \text{ for~}i \in [n]\\
  \label{eq:ax4}
  \left(1 - \sum_{j \le i} \isvar(v,a,j)\right)
  \left(\sum_{j < i} \isvar(v,a,j)\right) &= 0,
  \text{ for~}i \in [n]\\
  \label{eq:ax5}
  \left(1 - \sum_{j \le i} \isvar(v,a,j)\right)
  \isvar(v,a,i) &= 0,
  \text{ for~}i \in [n] \eqperiod
\end{align}
\Cref{eq:ax1} is equal to \cref{eq:ax-in-gt} and similarly
\cref{eq:ax2} is equal to \cref{eq:prod-zero-in}. In the following we
show that \cref{eq:ax3,eq:ax4,eq:ax5} can be derived from
\cref{eq:prod-zero-in}, the Boolean axioms and the negation axioms in
constant degree.

Consider \cref{eq:ax3}. Expand and rewrite modulo the Boolean axioms
and the negation axiom to obtain
% \begin{align}
%   \sum_{j \le i} \isvar(v,a,j) -
%   \isvar(v,a,i)\left(\sum_{j \le i} \isvar(v,a,j)\right) -\\
%   \left(\sum_{j \le i} \isvar(v,a,j)\right)
%   \left(\sum_{j < i} \isvar(v,a,j)\right) +\\
%   \isvar(v,a,i)
%   \left(\sum_{j \le i} \isvar(v,a,j)\right)
%   \left(\sum_{j < i} \isvar(v,a,j)\right)
%    = 0 \eqperiod
% \end{align}
% Rewrite modulo the Boolean axioms to get to
% \begin{align}
%   \sum_{j < i} \isvar(v,a,j) -
%   \isvar(v,a,i)\left(\sum_{j < i} \isvar(v,a,j)\right) -\\
%   \sum_{j < i} \isvar(v,a,j) -
%   2\sum_{j < j' < i}\isvar(v,a,j)\isvar(v,a,j') -
%   \isvar(v,a,i)\sum_{j < i}\isvar(v,a,j) +\\
%   \isvar(v,a,i)
%   \left(\sum_{j < i} \isvar(v,a,j)\right) +
%   \isvar(v,a,i)
%   \left(\sum_{j < i} \isvar(v,a,j)\right)^2
%   = 0 \eqperiod
% \end{align}
% We can cancel some terms and are left with the polynomial equation
\begin{align}
  &\isvar(v,a,i)
  \left(\sum_{j < i} \isvar(v,a,j)\right)^2-\nonumber\\
  &\qquad\qquad 2\sum_{j < j' < i}\isvar(v,a,j)\cdot\isvar(v,a,j')~-\nonumber\\
  &\qquad\qquad\qquad\qquad\isvar(v,a,i)\sum_{j < i}\isvar(v,a,j)
  = 0 \eqperiod
\end{align}
Observe that every term $t$ left in this polynomial is of the form
$t = t' \cdot \isvar(v,a,j) \cdot \isvar(v,a,j')$, for some
$j \neq j'\in [i]$ and a term $t'$ of degree at most $1$. But this
means that every term is equal to $0$ modulo \cref{eq:prod-zero-in}
and we thus see that \cref{eq:ax3} can be derived in constant degree
from $\circuit_s(f)$.

Let us consider \cref{eq:ax4}. Rewrite modulo the Boolean axiom to
obtain
\begin{align}
  &\isvar(v,a,i)\sum_{j < i}\isvar(v,a,j) +\nonumber\\
  &\qquad 2 \sum_{j < j' < i}\isvar(v,a,j)\cdot \isvar(v,a,j')
  = 0 \eqperiod
\end{align}
All terms are of the form of \cref{eq:prod-zero-in} and we can thus
derive \cref{eq:ax4} from $\circuit_s(f)$ in constant degree.

Last, we need to consider \cref{eq:ax5}. Note that modulo the Boolean
axiom we obtain the polynomial equation
\begin{align}
  - \isvar(v,a,i)\sum_{j < i}\isvar(v,a,j) = 0 \eqperiod
\end{align}
Also in this polynomial every term is of the form of
\cref{eq:prod-zero-in} and thus also \cref{eq:ax5} can be derived in
constant degree.

What remains is to show that \cref{ax:gate-cnf} can be derived from
$\circuit_s(f)$ in constant degree. This can be checked analogous to
\cref{ax:input-cnf} and we thus omit it here.

We conclude that all axioms of $\restrict{\circuit^{\mathrm{CNF}}_s(f)}{\rho}$ can
be derived from $\circuit_s(f)$ in constant degree and thus a degree
$d$ SoS refutation of $\circuit_s^{\mathrm{CNF}}(f)$ gives rise to a degree
$O(d)$ SoS refutation of $\circuit_s(f)$. Equivalently, a degree $d$
lower bound for $\circuit_s(f)$ implies a degree $\Omega(d)$ lower
bound for $\circuit_s^{\mathrm{CNF}}(f)$ as claimed.

\begin{proposition}
  Let $s \in N$ and $f: \set{0,1}^n \rightarrow \set{0,1}$ be a
  Boolean function on $n$ bits that requires circuits of size larger
  than $s$ to be computed. Then there is a degree $O(s)$ SoS
  refutation of $\circuit_s(f)$.
\end{proposition}

\section{Degree Upper Bound}
\label{sec:upper-bound}

In this section we give a simple upper bound on the SoS refutation
degree for $\circuit_s(f)$.  Specifically, for functions $f$ that have
no circuit of size $s$, we show that there is an SoS refutation of
$\circuit_s(f)$ of degree $O(s)$, essentially matching our
$\Omega(s^{1-\epsilon})$ lower bound.

At a high level, the logic behind the refutation is as follows: first,
we show that SoS in degree $O(s)$ can derive that the $\Theta(s^2)$
structure variables of $\circuit_s(f)$ must uniquely describe a
circuit of size $s$.  Then, we also show that for any fixed circuit,
SoS can in degree $O(s)$ derive that the circuit does not compute $f$
by ``evaluating'' the circuit on some (non-deterministically chosen)
input where the output differs from $f$.

To make this precise, we first define a set of monomials, which correspond
to circuits of size $s$. A multilinear monomial $m$ is a \emph{circuit
  monomial} if for every gate $v \in [s]$ it holds that
\begin{enumerate}
\item exactly one of the variables $\isneg(v), \isor(v)$ or $\isand(v)$ occurs
  in $m$,
\item for $a \in \set{1,2}$ exactly one of the variables
  $\isfromconst(v, a)$, $\isfromvar(v, a)$ or $\isfromgate(v, a)$ occurs
  in $m$,
\item for $a \in \set{1,2}$ exactly one of the variables
  $\constval(v, a)$ or $\overline{\constval(v, a)}$ occurs in $m$,
\item for $a \in \set{1,2}$ exactly one of the variables
  $\set{\isvar(v,a,i) \mid i \in [n]}$
  occurs in $m$, and
\item for $a \in \set{1,2}$ and $v > 1$, exactly one of the variables
  $\set{\isgate(v,a,u) \mid u < v}$
  occurs in $m$, and
  \item no other variables occur in $m$ than the ones described
    above.
\end{enumerate}
We denote by $\calM_s$ the set of circuit monomials. We first show
that SoS can derive in degree $O(s)$ the polynomial
$\sum_{m \in \calM_s} m - 1$; this corresponds to SoS proving that
the structure variables uniquely describe a circuit and does not use
anything about $f$.  Then in a second step we show that for every
$m \in \calM_s$, SoS can derive $-m$ in degree $O(s)$; this
corresponds to SoS proving that the circuit described by $m$ does not
compute $f$ correctly.  Summing these two parts up yields an SoS
derivation of $-1$, i.e., a refutation of the $\circuit_s(f)$ formula.

\pa{Here we still talk about SoS  derivations so we can't remove it from prelims}

\paragraph{Deriving~$\bm{\sum_{m \in \calM_s} m - 1}$.} We proceed
by induction on $s$. Note that $\calM_0 = \set{1}$ and hence the base
case is trivial. Suppose we have an SoS derivation of
$\sum_{m \in \calM_s} m - 1$. For every monomial $m \in \calM_s$ we
add the polynomial
\begin{align}
  m \cdot \big(\isneg(v) +  \isor(v) + \isand(v) - 1\big)
\end{align}
to the derivation (note that the second term is \cref{eq:ax-fn}). This
gives us an SoS derivation of $\sum_{m \in \calM'_s} m - 1$, where
$$\calM'_s = \bigcup_{m \in \calM_s} \set{m \cdot \isneg(v), m \cdot \isor(v), m \cdot \isand(v)}\eqperiod$$ We can continue in the same manner with
\cref{eq:ax-from} and \cref{eq:ax-in-gt} to finally obtain an SoS
derivation of $\sum_{m \in \calM_{s+1}} m - 1$. Clearly this
derivation requires degree at most $O(s)$, as for each gate there are
at most $7$ variables in every monomial from $\calM_s$.

\paragraph{Deriving $-m$ for $m \in \calM_s$.}
Let $C$ be the circuit that corresponds to the monomial $m$ and let
$\alpha \in \set{0,1}^n$ be such that $f(\alpha) \neq C(\alpha)$ (by
the assumption that $f$ does not have a circuit of size $s$, such an
$\alpha$ exists).  Suppose $C(\alpha) = b$ but $f(\alpha) = 1-b$.

We construct a degree $O(s)$ SoS proof of the fact that
$C(\alpha) = b$. That is, we are going to show that the polynomial
$p_s = m \cdot \left( \outwire_\alpha(s) - b\right)$ can be written as
\begin{align}
  \label{eq:-m-derivation}
  p_s = \sum_{i=1}^t r_i \cdot q_i \eqcomma
\end{align}
for some parameter $t$, axioms $q_1, \ldots, q_t$ and some polynomials
$r_1, \ldots, r_t$, such that $\deg(r_i \cdot q_i) = O(s)$ for all
$i$.  Note that given this, we can then easily derive $-m$ by
subtracting the polynomial
$m \cdot \left( \outwire_\alpha(s) - f(\alpha)\right) = m \cdot \left(
  \outwire_\alpha(s) - 1 + b\right)$ (this is $m$ multiplied by
\cref{eq:ax-correct}), yielding a derivation of $m \cdot (1 - 2b)$
which is either $m$ or $-m$ depending on $b$; in the former case it
can be multiplied by $-1$ to yield $-m$.  Note that here it is
important that the derivation \eqref{eq:-m-derivation} is a
Nullstellensatz derivation, not using any Sum-of-Squares part, since
otherwise it would not be possible to multiply it by $-1$.

Let us thus see how to derive $m \cdot \left( \outwire_\alpha(s) - b\right)$.
We do this by structural induction over the circuit: for every gate
$v$ we are going to construct an SoS proof of the fact that the
circuit rooted at $v$ outputs the bit $b_v$ on input $\alpha$. In
other words, an SoS derivation of the polynomial
$m \cdot \left(\outwire_\alpha(v) - b_v\right)$.


Let us explain how to construct an SoS proof $p_v$. Consider a gate
$v$ in the circuit. Depending on the function computed at $v$ and how
the wires of $v$ are connected we construct $p_v$ slightly differently.
As a first step let us construct SoS proofs $q_1$ and $q_2$ of the
fact that on input $\alpha$ the bit $c_a$, $a \in \set{1,2}$, is
carried on wire $a$ to the gate $v$. That is, the polynomial $q_a$
should simplify to $m \cdot \big(\inwire_\alpha(v, a) - c_a\big)$. In
the following we explain how to precisely define $q_a$ depending on
what the wire is connected to. Note that not a lot is going on -- we
are mostly just multilinearizing using the Boolean axioms.

If $m$ is of the form $m = m' \cdot \isfromconst(v,a) \cdot \constval(v, a)$ for some monomial $m'$, i.e., wire $a$
is connected to the constant $c_a = 1$ in the circuit described by $m$, then we
can derive $q_a = m \cdot \big(\inwire_{\alpha}(v,a) - 1\big)$
by the identity
\begin{align}
  \nonumber
  q_a &= m' \cdot \constval(v,a) \cdot \isfromconst(v,a) \cdot \big(\inwire_{\alpha}(v,a) -
        \constval(v, a)\big) +\\
      &\qquad m' \cdot \isfromconst(v,a) \cdot
        \big( \constval(v,a)^2 - \constval(v,a)\big) \eqcomma
\end{align}
a linear combination of \cref{eq:inwire-const} and the Boolean axiom on $\constval(v,a)$.  Similarly if $m = m' \cdot \isfromconst(v,a) \cdot \overline{\constval(v,a)}$ (i.e., $c_a = 0$) we can derive $q_a = m \cdot \inwire_{\alpha}(v,a)$ by
\begin{align}
  \nonumber
  q_a &= m' \cdot \overline{\constval(v,a)} \cdot \isfromconst(v,a) \cdot \big(\inwire_{\alpha}(v,a) -
        \constval(v, a)\big) +\\\nonumber
      &\qquad m' \cdot \isfromconst(v,a) \cdot \constval(v,a) \cdot
         \\\nonumber
      &\qquad\qquad\qquad \big( 1 - \constval(v,a) - \overline{\constval(v,a)}\big) +\\
      &\qquad m' \cdot \isfromconst(v,a) \cdot (\constval(v,a)^2 - \constval(v,a))
        \eqcomma
\end{align}
where we additionally use the negation axiom on $\constval(v,a)$.

Next, if $a$ is connected to an input $i$ and $m$ is of the form
$m = m' \cdot \isfromvar(v, a) \cdot \isvar(v, a, i)$ (so that $c_a = \alpha_i$), then
\begin{align}
  q_a = m \cdot (\inwire_{\alpha}(v,a) - \alpha_i) &= m' \cdot \isfromvar(v, a) \cdot \isvar(v, a, i) \cdot
          \big( \inwire_\alpha(v, a) - \alpha_i\big) \eqcomma
\end{align}
which is a multiple of \cref{eq:inwire-var}.  Lastly, if
$a$ is connected to a gate $u$ and $m$ is of the form
$m = m' \cdot \isfromgate(v, a) \cdot \isgate(v, a, u)$ (i.e., $c_a = b_u$), then
\begin{align}
  \nonumber
  q_a = m \cdot(\inwire_\alpha(v,a) - b_u) &= m' \cdot
          \isfromgate(v, a) \cdot \isgate(v, a, u) \cdot
          \big(\inwire_\alpha(v, a) - \outwire_\alpha(u)\big)\\
                                           &\qquad+ m \cdot (\outwire_\alpha(u) - b_u)
                                             \eqcomma
\end{align}
where the first term is a multiple of \cref{eq:inwire-gate}, and the
second term is the polynomial $p_u$ which, by induction, we assume has
already been derived in degree $O(s)$.

Given the two SoS proofs $q_1$ and $q_2$ we are ready to construct the
SoS proof $p_v = m \cdot (\outwire_\alpha(v) - b_v)$. As mentioned earlier
we do a case distinction over the funtion computed at $v$.
\begin{enumerate}
\item \emph{$v$ is a \emph{not} gate ($m = m' \cdot \isneg(v)$ and $b_v = 1-c_1$).} We have the derivation
  \begin{align}
    \nonumber
    p_v &=  m' \cdot \isneg(v) \cdot
             \big(\outwire_{\alpha}(v) - \overline{\inwire_{\alpha}(v, 1)}\big) \\\nonumber
           &\qquad + m \cdot
             \big(\overline{\inwire_{\alpha}(v, 1)} - 1 + \inwire_{\alpha}(v, 1)\big)\\
           &\qquad - q_1\eqcomma
  \end{align}
  where the first line uses \cref{eq:ax-neg} and the
  second line uses the negation axiom for $\inwire_{\alpha}(v, 1)$.
\item \emph{$v$ is an \emph{or} gate ($m = m' \cdot \isor(v)$, and $b_v = 1-(1-c_1)(1-c_2)$.} We have the derivation
  \begin{align}
    \nonumber
    p_v =  & m' \cdot \isor(v) \cdot 
             \Big(
             \outwire_{\alpha}(v) -
             \big(
             1 -
             \overline{\inwire_{\alpha}(v, 1)} \cdot
             \overline{\inwire_{\alpha}(v, 2)}
             \big)
             \Big)\\\nonumber
           & -
             m \cdot
             \overline{\inwire_\alpha(v, 1)} \cdot
             \big(
             \overline{\inwire_{\alpha}(v, 2)}
             - 1 + \inwire_{\alpha}(v, 2)
             \big)\\\nonumber
           & +
             m \cdot
             \big(
             \inwire_\alpha(v, 2) - 1
             \big)
             \cdot
             \big(
             \overline{\inwire_{\alpha}(v, 1)}
             - 1 + \inwire_{\alpha}(v, 1)
             \big)\\\nonumber
           & +
             \big(
             1 - \inwire_\alpha(v, 1)
             \big)
             \cdot
             q_2 \\
           & +
             \big(
             1 - c_2
             \big)
             \cdot
             q_1\eqcomma
  \end{align}
  where the first line uses \cref{eq:ax-or}, and the
  following two lines uses negation axioms.
\item \emph{$v$ is an \emph{and} gate ($m = m' \cdot \isand(v)$, and $b_v = c_1 \cdot c_2$.} We have
  \begin{align}
    \nonumber
    p_v =  & m' \cdot \isand(v) \cdot
             \big(
             \outwire_{\alpha}(v) -
             \inwire_{\alpha}(v, 1) \cdot
             \inwire_{\alpha}(v, 2)
             \big)\\\nonumber
           & +
             \inwire_\alpha(v, 1)
             \cdot
             q_2 \\
           & +
             c_2
             \cdot
             q_1\eqcomma
  \end{align}
  where the first line uses \cref{eq:ax-and}.
\end{enumerate}
This completes the description of the SoS derivation of
$\outwire_\alpha(s) = b$.
%
Observe that the final proof $p_s$ is of degree $O(s)$: in every
inductive step we increase the degree of the proof by at most a
constant.

\newcommand{\etalchar}[1]{$^{#1}$}
\begin{thebibliography}{KMOW17}

\bibitem[ABRW04]{ABRW04Pseudorandom}
Michael Alekhnovich, Eli {Ben-Sasson}, Alexander~A. Razborov, and Avi
  Wigderson.
\newblock Pseudorandom generators in propositional proof complexity.
\newblock {\em SIAM Journal on Computing}, 34(1):67\nobreakdash--88, 2004.
\newblock Preliminary version in \emph{FOCS~'00}.

\bibitem[AH19]{AH18SosTradeoff}
Albert Atserias and Tuomas Hakoniemi.
\newblock {Size-Degree Trade-Offs for Sums-of-Squares and Positivstellensatz
  Proofs}.
\newblock In Amir Shpilka, editor, {\em 34th Computational Complexity
  Conference (CCC 2019)}, volume 137 of {\em Leibniz International Proceedings
  in Informatics (LIPIcs)}, pages 24:1--24:20, Dagstuhl, Germany, 2019. Schloss
  Dagstuhl--Leibniz-Zentrum fuer Informatik.

\bibitem[AOW15]{AOW15}
S.~R. Allen, R.~ODonnell, and D.~Witmer.
\newblock How to refute a random csp.
\newblock In {\em 2015 IEEE 56th Annual Symposium on Foundations of Computer
  Science (FOCS)}, pages 689--708, Los Alamitos, CA, USA, oct 2015. IEEE
  Computer Society.

\bibitem[Ber82]{B82}
S.~J. Berkowitz.
\newblock On some relationships between monotone and non-monotone circuit
  complexity.
\newblock Technical report, Technical Report, University of Toronto, 1982.

\bibitem[BHK{\etalchar{+}}16]{BHKKMP16clique}
B.~{Barak}, S.~B. {Hopkins}, J.~{Kelner}, P.~{Kothari}, A.~{Moitra}, and
  A.~{Potechin}.
\newblock A nearly tight sum-of-squares lower bound for the planted clique
  problem.
\newblock In {\em 2016 IEEE 57th Annual Symposium on Foundations of Computer
  Science (FOCS)}, pages 428--437, 2016.

\bibitem[BT06]{BT06}
Andrej Bogdanov and Luca Trevisan.
\newblock Average-case complexity.
\newblock {\em Foundations and Trends in Theoretical Computer Science},
  2(1):1--106, 2006.

\bibitem[BW06]{BW06}
Amos Beimel and Enav Weinreb.
\newblock Monotone circuits for monotone weighted threshold functions.
\newblock {\em Inf. Process. Lett.}, 97(1):12–18, jan 2006.

\bibitem[GHP02]{GHA02proofs}
Dima Grigoriev, Edward~A. Hirsch, and Dmitrii~V. Pasechnik.
\newblock Complexity of semi-algebraic proofs.
\newblock In {\em STACS 2002}, pages 419--430, Berlin, Heidelberg, 2002.
  Springer Berlin Heidelberg.

\bibitem[Gol20]{Goldreich2020}
Oded Goldreich.
\newblock {\em On (Valiant's) Polynomial-Size Monotone Formula for Majority},
  pages 17--23.
\newblock Springer International Publishing, Cham, 2020.

\bibitem[Gri01]{gri01xor}
Dima Grigoriev.
\newblock Linear lower bound on degrees of positivstellensatz calculus proofs
  for the parity.
\newblock {\em Theoretical Computer Science}, 259(1):613 -- 622, 2001.

\bibitem[GUV09]{GUV09Unbalanced}
Venkatesan Guruswami, Christopher Umans, and Salil Vadhan.
\newblock Unbalanced expanders and randomness extractors from
  {P}arvaresh--{V}ardy codes.
\newblock {\em Journal of the ACM}, 56(4):20:1\nobreakdash--20:34, July 2009.
\newblock Preliminary version in \emph{CCC~'07}.

\bibitem[GW95]{GW95}
Michel~X. Goemans and David~P. Williamson.
\newblock Improved approximation algorithms for maximum cut and satisfiability
  problems using semidefinite programming.
\newblock {\em J. {ACM}}, 42(6):1115--1145, 1995.

\bibitem[Hir18]{Hi18}
Shuichi Hirahara.
\newblock Non-black-box worst-case to average-case reductions within np.
\newblock In {\em 2018 IEEE 59th Annual Symposium on Foundations of Computer
  Science (FOCS)}, pages 247--258, 2018.

\bibitem[KC00]{KC00}
Valentine Kabanets and Jin-Yi Cai.
\newblock Circuit minimization problem.
\newblock In {\em Proceedings of the Thirty-Second Annual ACM Symposium on
  Theory of Computing}, STOC '00, page 73–79, New York, NY, USA, 2000.
  Association for Computing Machinery.

\bibitem[KMOW17]{kmow17anycsp}
Pravesh~K. Kothari, Ryuhei Mori, Ryan O’Donnell, and David Witmer.
\newblock Sum of squares lower bounds for refuting any csp.
\newblock In {\em Proceedings of the 49th Annual ACM SIGACT Symposium on Theory
  of Computing}, STOC 2017, page 132–145, New York, NY, USA, 2017.
  Association for Computing Machinery.

\bibitem[KMS98]{KMS98}
David Karger, Rajeev Motwani, and Madhu Sudan.
\newblock Approximate graph coloring by semidefinite programming.
\newblock {\em J. ACM}, 45(2):246–265, mar 1998.

\bibitem[Kra01]{Krajicek01}
Jan Krajíček.
\newblock On the weak pigeonhole principle.
\newblock {\em Fundamenta Mathematicae}, 170(1-2):123--140, 2001.

\bibitem[MPW15]{MPW15SumOfSquaresPlantedClique}
Raghu Meka, Aaron Potechin, and Avi Wigderson.
\newblock Sum-of-squares lower bounds for planted clique.
\newblock In {\em Proceedings of the 47th Annual ACM Symposium on Theory of
  Computing ({STOC}~'15)}, pages 87\nobreakdash--96, June 2015.

\bibitem[MW15]{MW15}
Cody~D. Murrayand and R.~Ryan Williams.
\newblock {On the (Non) NP-Hardness of Computing Circuit Complexity}.
\newblock In David Zuckerman, editor, {\em 30th Conference on Computational
  Complexity (CCC 2015)}, volume~33 of {\em Leibniz International Proceedings
  in Informatics (LIPIcs)}, pages 365--380, Dagstuhl, Germany, 2015. Schloss
  Dagstuhl--Leibniz-Zentrum fuer Informatik.

\bibitem[Pot20]{potechin20}
Aaron Potechin.
\newblock {Sum of Squares Bounds for the Ordering Principle}.
\newblock In Shubhangi Saraf, editor, {\em 35th Computational Complexity
  Conference (CCC 2020)}, volume 169 of {\em Leibniz International Proceedings
  in Informatics (LIPIcs)}, pages 38:1--38:37, Dagstuhl, Germany, 2020. Schloss
  Dagstuhl--Leibniz-Zentrum f{\"u}r Informatik.

\bibitem[PR04]{PitRaz04}
Toniann Pitassi and Ran Raz.
\newblock Regular resolution lower bounds for the weak pigeonhole principle.
\newblock {\em Combinatorica}, 24(3):503--524, 2004.
\newblock Preliminary version in \emph{STOC~'01}.

\bibitem[Raz98]{Razborov98}
Alexander~A. Razborov.
\newblock Lower bounds for the polynomial calculus.
\newblock {\em Computational Complexity}, 7(4):291\nobreakdash--324, December
  1998.

\bibitem[Raz04a]{Raz04}
Ran Raz.
\newblock Resolution lower bounds for the weak pigeonhole principle.
\newblock {\em J. {ACM}}, 51(2):115--138, 2004.

\bibitem[Raz04b]{Razborov04ResolutionLowerBoundsPM}
Alexander~A. Razborov.
\newblock Resolution lower bounds for perfect matching principles.
\newblock {\em Journal of Computer and System Sciences},
  69(1):3\nobreakdash--27, August 2004.
\newblock Preliminary version in \emph{CCC~'02}.

\bibitem[Raz15]{Razborov15PRG}
Alexander~A. Razborov.
\newblock Pseudorandom generators hard for \mbox{$k$-{DNF}} resolution and
  polynomial calculus resolution.
\newblock {\em Annals of Mathematics}, 181(2):415\nobreakdash--472, March 2015.

\bibitem[Raz21]{razborov21youtube}
Alexander Razborov.
\newblock {P}, {NP} and {Proof Complexity}.
\newblock \url{https://youtu.be/ZVL_HsPC4xE?t=2646}, 2021.
\newblock Accessed April 2022.

\bibitem[Raz22]{razborov22website}
Alexander Razborov.
\newblock Open problems.
\newblock \url{https://people.cs.uchicago.edu/~razborov/teaching/index.html},
  2022.
\newblock Accessed April 2022.

\bibitem[RRS17]{RRS17}
Prasad Raghavendra, Satish Rao, and Tselil Schramm.
\newblock Strongly refuting random csps below the spectral threshold.
\newblock In {\em Proceedings of the 49th Annual ACM SIGACT Symposium on Theory
  of Computing}, STOC 2017, page 121–131, New York, NY, USA, 2017.
  Association for Computing Machinery.

\bibitem[RWY02]{RazbWigYao02}
Alexander~A. Razborov, Avi Wigderson, and Andrew~Chi{-}Chih Yao.
\newblock Read-once branching programs, rectangular proofs of the pigeonhole
  principle and the transversal calculus.
\newblock {\em Combinatorica}, 22(4):555--574, 2002.
\newblock Preliminary version in \emph{STOC~'97}.

\bibitem[Val84]{Valiant84}
Leslie~G. Valiant.
\newblock Short monotone formulae for the majority function.
\newblock {\em J. Algorithms}, 5:363--366, 1984.

\end{thebibliography}
\end{document}
